Agent roles in the Codynamic Book Machine are divided so that each responsibility has a clear owner, a narrow interface, and a predictable place in the workflow. The result is not a single monolithic authoring process but a coordinated set of specialized tasks that move a book from intent to outline, from outline to draft, and from draft to compiled artifact.

At the center of the system is the hypervisor. The hypervisor does not write prose itself; instead, it supervises the overall run, assigns work, and keeps the project moving through the correct phases. It is responsible for initiating tasks, routing work to the appropriate agent, and maintaining the high-level view of progress across the book. In practice, the hypervisor acts as the orchestration layer that ensures the rest of the system is working on the right section at the right time.

The outlining agent turns author intent and imported material into a structured book plan. Its job is to shape the canonical outline, organize chapters and sections, and preserve the dependencies that make later drafting possible. Where the hypervisor manages the run, the outlining agent manages the structure of the work itself. It establishes the sequence of topics and the relationships among them so that downstream agents can draft against a stable scaffold.

Section drafting agents are responsible for producing the LaTeX body of individual sections. Each section agent works within a bounded payload and converts outline intent, imported notes, and coordination feedback into prose. This division keeps drafting local and incremental: a section agent does not need to solve the whole book, only the section it owns. That constraint makes revision easier and allows the system to parallelize work across multiple sections when needed.

The gardener agent performs cleanup and consistency work. It responds to structural or stylistic issues that emerge after drafting, such as uneven transitions, duplicated ideas, missing connective tissue, or mismatches between neighboring sections. Rather than introducing new subject matter, the gardener refines what is already present so the manuscript reads as a coherent whole. In this sense, gardening is the maintenance layer of the writing process.

Diagram agents handle visual assets when a figure, chart, or schematic would materially improve understanding. Their role is to create or retrieve the media needed to support the text and to return a usable path and inclusion hint for the LaTeX source. This keeps visual production separate from prose drafting while still allowing the section to request a diagram when the argument benefits from one.

The design agent manages presentation-level concerns that affect the appearance and layout of the book. It is concerned with how the manuscript is rendered, how structural choices interact with the document format, and how the visual system supports readability. Where the diagram agent supplies content-specific visuals, the design agent helps ensure that the broader document design remains consistent.

The references agent owns citation registration and bibliography coordination. When a section needs a source that is not yet available in the shared reference set, the section workflow must request research and then submit candidate records to the references agent. This separation prevents ad hoc bibliography edits and keeps citation keys canonical across the project. The references agent therefore serves as the gatekeeper for source integrity and citation consistency.

Finally, the assembly agent compiles the finished manuscript into publication artifacts. It gathers the drafted sections, resolves the document structure, and produces the outputs used for review and release. Assembly is the point at which the distributed work of the other agents becomes a single book-shaped artifact.

Taken together, these roles form a pipeline with explicit boundaries: hypervision coordinates the run, outlining defines structure, section agents draft content, gardening improves coherence, diagram and design agents support presentation, the references agent controls source registration, and assembly produces the final output. The system is deliberately modular, but the modules are not independent. Each role exists to reduce ambiguity for the next role, so that the book can advance through a disciplined sequence of transformations rather than a loose collection of edits.

This division of labor also clarifies how the machine should be extended. New capabilities are added by assigning them to a role with a well-defined responsibility, rather than by expanding every agent's scope at once. That keeps the workflow legible: each agent knows what it owns, what it may request from others, and when its work is complete. In a long-form authoring system, that clarity is not just an implementation detail; it is what makes coordinated writing possible.
